\documentclass{article}
\usepackage[utf8]{inputenc}
\usepackage{amsmath}
\usepackage{amsfonts}
\usepackage{amssymb}
\usepackage{tikz}
\usepackage{fullpage}
\usepackage{tikz-cd}
\usepackage{spectralsequences}
\usepackage{adjustbox}
\usepackage{xfrac}
\usepackage{parskip}
\usepackage{tcolorbox}
\usepackage{xcolor}
\usepackage{hyperref}
\usepackage{graphicx}
\graphicspath{ {D:/Chrome Downloads./} }
\usepackage{amsthm}
\newcommand{\R}{\mathbb{R}}
\renewcommand{\P}{\mathbb{P}}
\newcommand{\Z}{\mathbb{Z}}
\newcommand{\C}{\mathbb{C}}
\newcommand{\A}{\mathbb{A}}
\newcommand{\Spec}{\textrm{Spec}}
\newcommand{\Proj}{\textrm{Proj }}
\newcommand{\Hom}{\textrm{Hom}}
\newcommand{\Aut}{\textrm{Aut}}
\newcommand{\Ext}{\textrm{Ext}}
\newcommand{\Tor}{\textrm{Tor}}
\newcommand{\Sym}{\textrm{Sym}}
\newcommand{\Der}{\textrm{Der}}
\newcommand{\End}{\textrm{End}}
\newcommand{\id}{\textrm{id}}
\usetikzlibrary{calc}
\theoremstyle{definition}
\newtheorem{theorem}{Theorem}[section]
\theoremstyle{definition}
\newtheorem{definition}{Definition}[theorem]
\theoremstyle{definition}
\newtheorem{remark}{Remark}[theorem]
\theoremstyle{definition}
\newtheorem{proposition}{Proposition}[theorem]
\theoremstyle{definition}
\newtheorem{conjecture}{Conjecture}[theorem]
\theoremstyle{definition}
\newtheorem{lemma}[theorem]{Lemma}
\theoremstyle{definition}
\newtheorem{corollary}{Corollary}[theorem]
\theoremstyle{definition}
\newtheorem{example}{Example}[theorem]
\tikzset{curve/.style={settings={#1},to path={(\tikztostart)
    .. controls ($(\tikztostart)!\pv{pos}!(\tikztotarget)!\pv{height}!270:(\tikztotarget)$)
    and ($(\tikztostart)!1-\pv{pos}!(\tikztotarget)!\pv{height}!270:(\tikztotarget)$)
    .. (\tikztotarget)\tikztonodes}},
    settings/.code={\tikzset{quiver/.cd,#1}
        \def\pv##1{\pgfkeysvalueof{/tikz/quiver/##1}}},
    quiver/.cd,pos/.initial=0.35,height/.initial=0}
\title{Rigidity of Cotangent Complex at Odd Primes}
\author{David Zhu}

\begin{document}
\maketitle


\tableofcontents

\newpage


\section{Introduction}
In his seminal work, Quillen proposed the following homology theory for augmented simplicial $R$-algebras, which is now dubbed Quillen homology: there is a Quillen adjunction 
\[
\operatorname{Ab} :
\mathrm{sCAlg}_R^{\mathrm{aug}}
\rightleftarrows
\mathrm{sMod}_R :
R \oplus (-)
\]
For an object $X\in \mathrm{sCAlg}_R^{\mathrm{aug}}$, the Quillen homology of $X$ is defined to be the module homology of $\mathbb{L}\textrm{Ab}(X)$. With this construction, Quillen gave the definition of the cotangent complex $\mathbb{L}_{A/R}$ associated to morphism of commutative rings $R\to A$. let $P\to A$ be a cofibrant replacement of $A$ in $s \textbf{CAlg}_R$. The cotangent complex $\mathbb{L}_{A/R}$ is modelled by the simplicial $A$-module
\[\mathbb{L}_{A/R}:= \Omega_{P/R}\otimes_{P}A\]
By the Dold-Kan correspondence, we will use $\mathbb{L}_{A/R}$ to denote both the simplicial $A$-module and its normalization in $\textrm{Ch}^+(A)$.The Quillen homology is defined to be
\[H_i^Q(A/R;k):= \pi_i(\mathbb{L}_{A/R}\otimes_{A}k)\cong \Tor^A_i(\mathbb{L}_{A/R};k)\]
In his seminal work, Quillen also proposed the following rigidity conjecture,

\begin{conjecture}
If $\textrm{pd}(\mathbb{L}_{A/R})<\infty$, then $\textrm{pd}(\mathbb{L}_{A/R})\leq 2$.
\end{conjecture}





\section{Outline}
Fix a morphism of Noetherian local $k$-algebras $R\to A$. Standard localization techniques reduces the conjecture to the case where we consider a surjection of Noetherian local $k$-algebras $(R,\mathfrak{m})\to (A, \mathfrak{n})$. In this case, the flat dimension of the cotangent complex can be detected by tensoring with the residue field
\begin{proposition}[Stack 15.79.2]
    Let $R$ be a ring. Let $a,b\in \Z$. Let $K^{\bullet}$ be a pseudo-coherent complex of $R$-modules. The following are equivalent: 
    \begin{enumerate}
        \item $K^{\bullet}$ is perfect with Tor amplitude in $[a,b]$.
        \item For every maximal ideal $\mathfrak{m}$, we have $\Tor^R_i(K^{\bullet};k(\mathfrak{m}))=0$ for $i\not \in [a,b]$
    \end{enumerate}
\end{proposition}
Thus, under perfectness assumption, it suffices to show $H_i^Q(A/R;k)=0$ for $i>2$. We will therefore suppress the coefficient $k$ in Quillen homology from now on, unless other coefficients are used.


To prove the conjecture, Turner's program utilizes many tools from the homotopy theory of simplicial commutative algebras. For a $k$-vector space $V$, let $K(V,n)$ be the simplicial $k$-vector space characterized by having homotopy group concentrated in degree $n$, with value $V$. 

\begin{theorem}
    Given $A\in s \textbf{CAlg}_k$, which is connected, there exists a sequence of cofibrations 
    \[A=A(0)\xrightarrow{j_0} A(1)\xrightarrow{j_1} A(2)\xrightarrow{j_2}...\]
    with the following properties
    \begin{enumerate}
        \item For each $n\geq 1$, $A(n)$ is $(n-1)$-connected.
        \item For $s\geq n$, 
         \[H_s^Q(A)\cong H_s^Q(A(n))\]
        \item There is a cofiber sequence 
        \[ \textrm{Sym}_k(K(H_n^Q(A),n))\to A(n)\to A(n+1) \]
    \end{enumerate}
\end{theorem}
In particular, for $A$ such that $H_i^Q(A)$ vanishes for $i>n$, we may identify 
\[A(n)\cong \textrm{Sym}_k(K(H_n^Q(A),n))\]
and obtain a comparison map 
\[\phi_n: A\to \textrm{Sym}_k(K(H_n^Q(A),n))\]
inducing isomorphism on $H_n^Q$. The work of Bousfield, Goerss and Lada identifies the homotopy groups of the right hand side, which is generated by compositions of homotopy operations.

\begin{theorem}[Bousfield, 8.11]
    Let $\{a\}_{a\in A}$ be $\mathbb{F}_p$-basis for $\pi_*X$. Then $\oplus_{s=0}^{\infty} \pi_t(\Sym^s X)$ is the non-unital associative graded-commutative $\mathbb{F}_p$-algebra with generators 
    \[\{\alpha_{I}(a) | a\in A, I \textrm{ admissible}, \textrm{ excess}(I)\leq |a| \}\]
\end{theorem}
In particular, one could find a composition of homotopy operations, of arbitrary length, acting non-trivially on the degree $n$-element in $\pi_*(\textrm{Sym}_k(K(H_n^Q(A),n)))$. Now the proof reduces to two parts:

\begin{theorem}
    Suppose $n>2$. There is non-zero element in $y\in \pi_* A$ such that $(\phi_n)_*(y)=a$, and one could find a composition of homotopy operations, of arbitrary length, acting non-trivially on $y$.
\end{theorem}
The proof of the theorem is comparing the Quillen spectral sequences using $\phi_n$.
\begin{theorem}
    There is a convergent spectral sequence
    \[E^1_{s,t}:= \pi_t(\Sym^s ( H^Q_*(A/R;k)) )\Rightarrow \pi_{s+t} X\]
\end{theorem}


This is the step where Turner had to restrict to characteristic $2$: he needed explicit Adem relations, which is still absent to this day, among operations to control the differentials. 


On the other hand, for $A= P\otimes_R k$, operations on $\pi_* A$ has the following nilpotence property
\begin{theorem}
  There exists some integer $N$ such that any composition of strictly suspended operation of length $>N$ on $\pi_* A$ is zero. 
\end{theorem}
This is where Turner needed the assumption that $A$ is Cohen-Macauley, so he could reduce to the Artinian case where the maximal ideal is nilpotent.




\section{\texorpdfstring{Homotopy Operations on Simplicia $\mathbb{F}_p$-Algebras }{} }

\begin{theorem}[Bousfield, 8.11]
    For an odd prime $p$, and a simplicial $\mathbb{F}_p$-algebra $B$, there exists two families of operations 
    \[v_t: \pi_n B\to \pi_{n+2t(p-1)}B \quad \quad 2\leq 2t\leq n \]
    \[\eta_t: \pi_n B\to \pi_{n+2t(p-1)+1}B \quad \quad 3\leq 2t+1\leq n \]
The operation $v_t$ are additive for $t< n$, and the top operation $v_n$ is the divided power.



\end{theorem}

\begin{definition}
    For $k\geq 0$, let $I=(\delta_1,s_1,...,\delta_k,s_k)$ be a $(2k)$-tuple of integers, where $\delta_i =0,1$ and $s_i\geq 1$. The index $I$ is called \textbf{admissible} if $s_{i}\geq ps_{i+1}+ \delta_i$ for $1\leq i\leq k-1$.  Its \textbf{excess} is defined to be 
    \[\textrm{excess}(I):=\begin{cases}
        0 & k=0 \\
        2s_1+\delta_1 & k=1 \\
        2s_1+\delta_1 - \sum^k_{i=2}(2(p-1)s_i+\delta_i) & k\geq 2
    \end{cases}
    \]
    
    
    
    
    Define 
    \[\alpha_{\delta_i,s_i}:=\begin{cases}
        v_{s_i} & \delta_i =0\\
        \eta_{s_i} & \delta_i=1
    \end{cases}\]
    and 
    \[\alpha_I:= \alpha_{\delta_1,s_i}\circ ...\circ \alpha_{\delta_k,s_k} \]

    \textcolor{red}{If out of range just zero}

\end{definition}


\begin{theorem}[Bousfield, 8.11]
    \label{generator}
    Let $\{a\}_{a\in A}$ be $\mathbb{F}_p$-basis for $\pi_*X$. Then $\oplus_{s=0}^{\infty} \pi_t(\Sym^s X)$ is the non-unital associative graded-commutative $\mathbb{F}_p$-algebra with generators 
    \[\{\alpha_{I}(a) | a\in A, I \textrm{ admissible}, \textrm{ excess}(I)\leq |a| \}\]
\end{theorem}

In particular, the generators provide an infinite tower of non-zero operations.
\begin{corollary}
    \label{tower}
    Let $V$ be a $k$-vector space, and let $u$ be the generating class in $\pi_n(\Sym_k K(V,n))$. For any $j>1$, the index $I_j:=(0,p^{j},0,p^{j-1},....,0,p,0,1)$ is admissible, and thus the operation $\theta_j:=\alpha_{I_j}$ is non-zero on $u$ for all $j$.
\end{corollary}

\begin{lemma}[Degree-Weight Bound]
    \label{degree_weight}
   Let $X$ be a simplicial $k$-module. 
   \begin{enumerate}
    \item Suppose $d=\textrm{dim } \pi_1(X)<\infty$. Then $\pi_t(\Sym^s X)$ is zero if $t<2s-d$.
    \item Suppose $X= K(V,n)$. Then $\pi_t(\Sym^s X)$ is zero if $t<2s+n-2$.
   \end{enumerate}
   
   
   
\end{lemma}
\begin{proof}
    Let $y$ be a non-zero monomial in $\pi_*(\Sym_k X)$. By Theorem \ref{generator}, we can express $y$ as product of $k$ admissible monomial generators
    \[y=\alpha_{I_1}(x_1)...\alpha_{I_k}(x_k)\]
    where each $I_j$ is admissible. Let $I_j:=(\delta_1,s_1,...,\delta_n,s_n)$. Admissiblity provides a lower bound on degree in terms of the weight 
      \begin{equation*}
    \begin{aligned}
        \textrm{deg}(\alpha_{I_i}(x_i)) & = \textrm{deg}(x_i)+2(p-1)\sum_{i=0}^{n-1}s_i \\
        &\geq \textrm{deg}(x_i)+2(p-1)\sum_{i=1}^{n-1}p^i \\
        &=\textrm{deg}(x_i) + 2(p^{n}-1) \\
        &=\textrm{deg}(x_i)+ 2 \textrm{wt}(\alpha_I) -2 
    \end{aligned}
\end{equation*}

Since degree and weight both add in the product, we get 
\[\textrm{deg}(y)\geq 2 \textrm{wt}(y) + \sum^k_{i=1} \textrm{deg}(x_i)-2 \]
If $X=K(V,n)$, then $|x_i|=n$ for all $i$, and part $(2)$ follows easily.


If $d=\textrm{dim}(H^Q_1(Q))<\infty$. We see only $|x_i|=1$ will contribute to a negative summand, and each basis vector can only appear once in the product for it to be non-zero by graded commutativity. Therefore
\[\textrm{deg}(y)>2 \textrm{wt}(y)-d\]
so part $(1)$ follows.




  
\end{proof}














 

\section{Detection in Quillen Spectral Sequence}
Let $A$ be an augmented simplicial $\mathbb{F}_p$-algebra. The Quillen spectral sequence reads 
\[E^1_{s,t}:= \pi_t(\Sym^s ( H^Q_*(A)) )\Rightarrow \pi_{s+t} X\]
Suppose $n$ is the top non-vanishing degree of $H^Q(A)$, and let $u\in H_n^Q(A)$ be the generating class. Then, the element $\theta_j$ defined in Corollary \ref{tower} is an element in bidegree $(s,t):=(p^j,n+2(p^j-1))$. We can analyze the permanence property of this cycle: 

\begin{proposition}
    The element $\theta_j$ can be represented by a cycle $a\in F^{s}N(P)_t$, and $da\in F^{2s}N(P)_{t-1}$.
\end{proposition}
In particular, the spectral sequence differential $d_r$ is trivial on the representing class $a$ for $r<s$, as the target $F^{2s}N(P)_{t-1}\subseteq F^{s+r+1}N(P)_{t-1}$ is trivial in the subquotient $\frac{F^{s+r}N(P)_{t-1}}{F^{s+r+1}N(P)_{t-1}}$. 

\begin{proposition}
    The Quillen spectral sequence of $\Sym K(H^Q_{n}(A),n)$ degenerates on $E^1$. 
\end{proposition}

For $r\geq s$, the target groups will be a subquotient of 
\[E^1_{s+r,2s+n-3}\]
By Lemma \ref{degree_weight}, we see the target groups automatically vanish once we choose $j$ large enough. 

On the other hand, the there are no incoming differentials before page $s$, by comparing the spectral sequence with that of $\Sym K(H^Q_{n}(A),n)$. There are no incoming differentials on or after page $s$ simply by degree reasons. 

\textbf{extension}: by comparing the spectral sequences, we see that there is no extension problems: the difference lands in $\pi_{n+2s-2} \Sym^{s+1}(K(H_n))$. This group is zero by part $(2)$ of Lemma \ref{degree_weight}. 


\section{Nilpotence of Operations on Local Derived Fiber }
Let $(R,\mathfrak{m})$ be Noetherian local ring, $M$ a finitely generated R-module. Choose a resolution of $M$ be finite free $R$-modules 
\[...\xrightarrow{d_2}F_1\xrightarrow{d_1}F_0\xrightarrow{d_0}M\to 0\]




\begin{theorem}[Uniform Artin–Rees for Syzygies]
For a fixed ideal $J\subseteq R$, there is a uniform integer $h\geq 0$ such that for every $n\geq h$, and $i\geq 0$, one has 
    \[J^nF_i\cap \textrm{im}(d_{i+1})\subseteq J^{n-h}\textrm{im}(d_{i+1})\]
\end{theorem}

\begin{corollary}
    There is a uniform integer $h\geq 0$, such that for every $q\geq h$ and $i\geq 0$, the reduction
    \[\Tor_i^R(M;R/\mathfrak{m}^q)\to \Tor_i^R(M;k) \]
    is zero.
\end{corollary}








\appendix

\section{Quillen Spectral sequence}
By the standard model structure on $s \textbf{CAlg}_R$, one may choose a cofibrant replacement $P\to A$ such that $P$ is semi-free, i.e every $P_n$ is a free $R$-algebra.  Let $X:= P\otimes_R k$. Denote by $I$ the augmentation ideal and $QX:=I/I^2$ the indecomposables. Standard computation gives
\[QX\cong \Omega_{P/R}\otimes_P k\]
And therefore we have $H_i^Q(A/R)\cong \pi_i(QX)$. This provides the interpretation that Quillen homology of an augmented simplicial $k$-algebra $X$ is the homotopy group of its indecomposables $QX$. we have a descending filtration of $X$ by 
\[X\supseteq I\supseteq I^2\supseteq I^3\supseteq ...\]
which leads to a spectral sequence 
\[E^1_{s,t}:= \pi_t(I^s/I^{s+1})\Rightarrow \pi_{s+t} X\]
A standard fact on the model structure on $s \textbf{CAlg}_R$ says one may choose $P$ to be semi-free, meaning each simplicial degree $P_n$ is a free $R$-algebra. It follows that 
\[I^s/I^{s+1}\cong \Sym_k^s(I/I^2) \]
Since a chain complex of vector spaces splits as direct sums of its homology, $I/I^2$ also splits as the direct sum of its homotopy groups under the Dold-Kan correspondence.




\end{document}